%! root=../main.tex
\section{\pname Implementation}\label{sec:impl}

\comment{%
    * We can tightly connect the text here to workflow/architecture/execution model.

        - framework (Picinae)
        - implementation (armv7.v, arm7cfi.v --- Section #)
        - proofs of *implementation* correctness (armv7cfi_proof.v --- Section #)
        --> extraction
        - *.ml files + py glues --> implement binary rewriter (rewriter.sh)
            - **py glues** are outside the scope of the research
                - the proof will be about the program loader would load the generated ELF execution and initialize the process initial state to be compliant to assumptions associated with *the proof of the correctness.*
                    - what is *the proof of the correctness* ; ~~Picinae~~or ARMCFI?
            - Disassembler part that is implemented in Coq hasn't yet been proven

    * As discussed in \autoref{sec:overview}, this work comprises two major components: \nm{1} static binary CFI instrumentation, and \nm{2} proof artifacts that verify the conformance of the CFI binary rewriter to its safety specification.

    * a typical \pname workflow first verifies the correctness of \pname's components through the extraction process with the assistance of the Coq proof assistant, after which the extracted OCaml source is compiled and linked into the rewriter executable.
        - To be verified, the CFI binary rewriter is first lifted into Picinae's intermediate representation (PIR) and expressed as a function that takes symbolic variables as input and produces instrumented output.
        - Internally, the rewriter first decodes the input byte stream (i.e., the \texttt{.text} section) to identify target instructions, i.e., the locations to be instrumented.
}

As discussed in \autoref{sec:overview}, \pname comprises two major components:
(1) a static binary CFI rewriter that instruments branch instructions at the binary level, and
(2) proof artifacts that machine-check the conformance of the rewriter to its safety specification.
Both components are developed within the \Picinae{} framework~(\autoref{sec:picinae}), which provides the underlying infrastructure for lifting binary code into a Rocq-verifiable intermediate representation.

\subsection{Rocq Artifacts}

The verified core of \pname is implemented in Gallina, the specification language of the Rocq proof assistant.
The implementation is organized into three primary artifacts.

\heading{ARMv7 Semantics (\texttt{decoder.v}, \texttt{lifter.v}).}
The first artifact provides a formal semantics of the ARMv7 instruction set expressed in Picinae's \IL.
% Each instruction is specified as a Gallina function that maps a machine state (register file and memory) to an updated state, encoding all architectural effects including register writes, memory operations, and control-flow transfers.
The behavior of each instruction is specified as a Gallina function that produces \IL based on the instruction operands, encoding all architectural effects including register writes, memory operations, and control-flow transfers.
% This semantics is shared across the verification of both the decoder and the rewriter, forming the trusted base for all subsequent proofs. ???

\heading{CFI Rewriter (\texttt{cfi.v}).}
The second artifact implements the CFI binary rewriter described in \autoref{sec:overview} as a Gallina function.
Given the original instruction stream, base addresses $(\bi, \bip, \tbi)$, and a CFI policy $\pol$, the function produces the rewritten code blocks $\txtp$ and policy tables $\tables$.
Internally, the rewriter decodes the input byte stream to identify each instruction and classify it into one of the cases described in \autoref{sec:overview}: static control flow, PC-relative data access, dynamic control flow, or unsupported instruction.
It then constructs the corresponding instrumented block and synthesizes the hash table entries for dynamic branches.
\kjee{Note: the disassembler component embedded here has not yet been formally verified; its correctness is assumed as a pre-condition of the overall proof.}

\heading{Correctness Proofs (\texttt{cfi\_proof.v}).}
The third artifact contains the machine-checked proofs establishing that the rewriter satisfies the CFI safety theorem of \autoref{sec:formalization}.
Concretely, for any input program and any execution trace of the rewritten binary, the proofs guarantee that every control-flow transfer is permitted by $\pol$.
These proofs are discharged entirely within Rocq, providing a machine-verified certificate that ships alongside the rewriter.

\subsection{Extraction and the Binary Rewriter}

A typical \pname build first checks all proofs in \texttt{cfi\_proof.v} using the Rocq proof checker, then invokes Rocq's extraction mechanism to translate the verified Gallina functions into executable OCaml code.
The extracted OCaml modules are compiled and linked with a small set of unverified OCaml and Python driver scripts that handle the ELF-specific tasks outside the scope of the formal verification:
parsing the input ELF binary to locate the \texttt{.text} segment and its load address,
invoking the extracted \texttt{cfi\_rw} function with the appropriate arguments,
and assembling a new ELF binary that splices in the rewritten text and dispatch table segments.

The formal guarantees of \pname apply to the verified core.
The surrounding driver scripts are trusted components, analogous to the program loader:
the safety proof assumes that the loader initializes the process's initial memory state in a manner consistent with the memory layout assumptions of the CFI theorem (see \autoref{sec:formalization}).
Verifying these loader-level assumptions is left as future work.
